% Continuous-time (delta STRAIN, filtered assignments) vs play-level blocker metrics, all vs PFF.
\begin{table}[h!]
\centering
\footnotesize
\begin{tabular}{lccccc}
\toprule
 & \multicolumn{2}{c}{vs pressures allowed} & \multicolumn{2}{c}{vs beaten} & within-pos. \\
\cmidrule(lr){2-3}\cmidrule(lr){4-5}
Blocker rating & Pearson & Spearman & Pearson & Spearman & press. (P) \\
\midrule
Continuous dose ($\theta_t$) & -0.347 & -0.331 & -0.412 & -0.434 & -0.245 \\
Realized impedance (FN) & -0.227 & -0.226 & -0.136 & -0.042 & -0.423 \\
Realized impedance & -0.129 & -0.118 & -0.128 & -0.043 & -0.290 \\
Play-level plus-minus & -0.143 & -0.105 & -0.052 & -0.019 & -0.277 \\
\bottomrule
\end{tabular}
\caption{Blocker ratings against independent PFF charting on a common set of 207 linemen ($\geq 50$ snaps). All ratings are oriented higher${=}$better, so a correct association is NEGATIVE. The continuous-time \emph{dose} rating --- the blocking dose $\theta(b,j,t)$ a rusher receives, conditioning on the time-$t$ state (the sequentially-ignorable estimand of the causal connection) --- is the strongest overall predictor of both \emph{pressures allowed} and the instantaneous \emph{beaten} event, and it separates protection skill \emph{within} position (last column, correctly signed) where the first-difference $\Delta\theta$ specification inverts. The front-normalized realized impedance remains the sharpest \emph{within}-position discriminator. The dose rating correlates positively with the play-level plus-minus ($+0.35$), so the two broadly agree on overall blocking quality --- now recovered from frame-level strain dynamics rather than play aggregates.}
\label{tab:blocker_cont_vs_play}
\end{table}
